\documentclass[11pt,a4paper]{report}

% =====================================================
% Pacotes
% =====================================================
\usepackage[utf8]{inputenc}
\usepackage[T1]{fontenc}
\usepackage[portuguese]{babel}
\usepackage{geometry}
\geometry{left=3cm,right=3cm,top=3cm,bottom=3cm}
\usepackage{graphicx}
\usepackage{float}
\usepackage{caption}
\usepackage{subcaption}
\usepackage{listings}
\usepackage{xcolor}
\usepackage{hyperref}
\usepackage{tikz}
\usetikzlibrary{trees}

% =====================================================
% Bibliografia (biblatex + biber)
% =====================================================
\usepackage{biblatex}
\addbibresource{referencias.bib} % ficheiro .bib

% Espaçamento entre linhas
\linespread{1.1}

% =====================================================
% Código
% =====================================================
\lstset{
    basicstyle=\ttfamily\small,
    breaklines=true,
    frame=single,
    numbers=left,
    numberstyle=\tiny,
    showstringspaces=false
}

\hypersetup{
    colorlinks=true,
    linkcolor=blue,
    urlcolor=cyan
}

% =====================================================
% Documento
% =====================================================
\title{
\textbf{Dashboard Interativo para Análise de Débitos Diretos}\\
\large Laboratórios de Informática -- LESI 2025/26
}

\author{
Henrique Pereira (35009)\\
Rodrigo Correia (35008)\\
António Torres (35002)\\
Dinis Moreira (36742)
}

\date{}

\begin{document}

\maketitle
\setcounter{page}{2}

% =========================================
% Página 3: juntar todas as listas numa só página
% =========================================
\begingroup
\let\clearpage\relax  % impede quebras de página automáticas temporariamente
\tableofcontents
\listoffigures
\listoftables
\lstlistoflistings
\endgroup  % restaura \clearpage automaticamente

% =====================================================
\chapter{Introdução}
% =====================================================
Este relatório apresenta o desenvolvimento de um dashboard interativo para análise de débitos diretos, realizado no âmbito da unidade curricular de Laboratórios de Informática. O projeto teve como objetivo principal a leitura, validação e visualização de dados provenientes de ficheiros no formato PS2, permitindo uma análise clara e interativa da informação financeira.

A aplicação foi desenvolvida em Python e Shiny, utilizando princípios do paradigma funcional e boas práticas de desenvolvimento colaborativo \cite{python,shiny}. Para algumas tarefas de programação assistida e sugestões de código, recorreu-se a ferramentas de Inteligência Artificial como ChatGPT e GitHub Copilot \cite{chatgpt,copilot}. O trabalho apoiou-se nos slides fornecidos pelo professor \cite{slides}.

\section{Objetivos}
Os principais objetivos deste trabalho são:
\begin{itemize}
    \item Desenvolver um dashboard interativo para análise de débitos diretos;
    \item Processar e validar ficheiros no formato PS2;
    \item Aplicar o paradigma funcional/declarativo;
    \item Utilizar boas práticas de desenvolvimento colaborativo;
    \item Produzir documentação técnica clara e estruturada em LaTeX \cite{lamport1994latex}.
\end{itemize}

\newpage

% =====================================================
\chapter{Análise do Problema}
% =====================================================
\section{Formato PS2}
O formato PS2 é um padrão utilizado na comunicação entre empresas e instituições bancárias, usando registos de largura fixa, comumente aplicado em operações de débito direto.

\section{Estrutura dos Registos}
Os ficheiros PS2 contêm três tipos principais de registos:
\begin{description}
    \item[Tipo 1 -- Cabeçalho] Contém informações gerais sobre o ficheiro;
    \item[Tipo 2 -- Detalhe] Representa cada operação individual;
    \item[Tipo 9 -- Fim/Rodapé] Permite o fecho do ficheiro.
\end{description}

\section{Exemplo de Ficheiro PS2}
\begin{lstlisting}[caption={Exemplo de ficheiro PS2}]
120251026IPCA Energy 50123456700000000045375000002
2000000100100350651000025874125423568914200000000024750
2000000100200330000452108745961550184592300000000020625
900000000045375000002
\end{lstlisting}

\newpage

% =====================================================
\chapter{Desenvolvimento}
% =====================================================
\section{Arquitetura da Solução}
A solução foi desenvolvida com uma arquitetura modular.

\begin{figure}[H]
\centering
\begin{tikzpicture}[
level distance=1.5cm,
level 1/.style={sibling distance=4cm}
]
\node {Dashboard}
child {node {Leitura}}
child {node {Validação}}
child {node {Visualização}};
\end{tikzpicture}
\caption{Arquitetura modular da solução}
\end{figure}

\section{Processamento de Ficheiros}
O processamento dos ficheiros PS2 foi implementado com funções puras.

\begin{lstlisting}[language=Python, caption={Processamento de ficheiro PS2}]
def processar_ficheiro_ps2(caminho):
    with open(caminho, 'r') as f:
        linhas = f.readlines()
    return linhas
\end{lstlisting}

\section{Validação de Dados}
A validação assegura que o ficheiro PS2 apresenta uma estrutura adequada e que os totais coincidam corretamente.

\section{Tecnologias Utilizadas}
\begin{table}[H]
\centering
\caption{Tecnologias utilizadas}
\begin{tabular}{|l|l|}
\hline
\textbf{Tecnologia} & \textbf{Finalidade} \\ \hline
Python & Processamento de dados \\ \hline
Shiny & Dashboard interativo \\ \hline
Git/GitHub & Controlo de versões \\ \hline
LaTeX & Documentação \\ \hline
\end{tabular}
\end{table}


\section{Tarefas Realizadas}
\textbf{Henrique Pereira}: Como líder do grupo, participei ativamente na organização e estrutura do projeto desde o início. Trabalhei no dashboard desde o início ao fim, passando por diversas fases e versões. Durante o desenvolvimento estive presente em diversas funções como a implementação de soluções em Python, documentação em Doxygen, relatório codificado em Latex, entre outros. Na reta final garanti que o README.md estava devidamente atualizado e que todas as implementações estavam funcionais. \\
\textbf{Rodrigo Correia}: Contribuí para um bocado de tudo mas foquei me principalmente no desenvolvimento do relatório, troubleshooting, ideias de implementação e/ou otimização e desenvolvimento do dashboard. Na reta final decidi ficar mais focado na finalização do relatório já que estava atrasado. \\
\textbf{António Torres}: Durante o projeto, fui um dos  responsáveis pela criação e desenvolvimento do Dashboard. Contribuí para o desenvolvimento das funções de leitura e validação. Na fase final, devido à indisponibilidade temporária do meu computador, apoiei a equipa através de pesquisas externas que permitiram otimizar o código e melhorar o relatório \\
\textbf{Dinis Moreira}: Ao longo do projeto estive envolvido principalmente em tarefas de pesquisa, troubleshooting e validação do código. Contribuí para o desenvolvimento da função de leitura e implementei funções de validação para garantir a correção dos dados. Participei ativamente na elaboração e revisão do relatório. Na fase final, aprofundei a documentação com doxygen e dei apoio com ideias e melhorias para o dashboard. \\

\section{Dificuldades}
\textbf{Henrique Pereira}: Utilização adequada do Doxygen \\
\textbf{Rodrigo Correia}: Dificuldade em encontrar informação relativamente ao syntax/erros do LaTeX \\
\textbf{António Torres}: Desafios relacionados com a implementação de código em Python \\
\textbf{Dinis Moreira}: Utilização adequada do Doxygen, implementação de código na biblioteca Shiny \\
\newpage

% =====================================================
\chapter{Resultados e Conclusão}
% =====================================================
O dashboard desenvolvido permite a visualização interativa dos débitos diretos.

\begin{figure}[H]
\centering
\includegraphics[width=0.8\textwidth]{/home/vboxuser/Desktop/dashboard_final.png} % Substitua pelo caminho da tua imagem
\caption{Interface do dashboard desenvolvido}
\end{figure}

O trabalho atingiu os objetivos definidos, consolidando conhecimentos de programação em Python, desenvolvimento de dashboards interativos e trabalho colaborativo com o GitHub. Para trabalhos futuros, podem ser implementadas funcionalidades adicionais, como suporte a novos formatos de ficheiros e análises preditivas.

\newpage

% =====================================================
% Bibliografia
% =====================================================
\printbibliography[heading=bibintoc,title={Bibliografia}]

\end{document}
